% Part: normal-modal-logic
% Chapter: syntax-and-semantics
% Section: entailment

\documentclass[../../../include/open-logic-section]{subfiles}

\begin{document}

\olfileid{nml}{syn}{ent}

\olsection{तार्किक निष्पन्नता}

\begin{explain}
  एखाद्या जगात सत्य असण्याच्या व्याख्येने !!{formula}s मधील
  तार्किक निष्पन्नतेचा संबंध परिभाषित करता येतो.
  !!^a{formula}~$!B$ पासून~$!A$ तार्किकरीत्या निष्पन्न
  होते तेव्हा आणि केवळ तेव्हाच, जेव्हा $!B$ सत्य असेल
  तेव्हा $!A$ देखील सत्य असते. येथे ``तेव्हा'' म्हणताना
  कोणतेही प्रतिमान आणि त्यातील कोणतेही जग अभिप्रेत आहे.
\end{explain}

\begin{defn}
  $\Gamma$ हा !!{formula}s चा संच आणि $!A$ हे !!a{formula}
  असेल, तर $\Gamma$ पासून~$!A$ \emph{तार्किकरीत्या
  निष्पन्न होते}, म्हणजे $\Gamma \Entails !A$, तेव्हा आणि
  केवळ तेव्हाच, जेव्हा प्रत्येक प्रतिमान
  $\mModel{M} = \tuple{W, R, V}$ आणि प्रत्येक $w \in W$
  साठी, सर्व $!B \in \Gamma$ बाबत $\mSat{M}{!B}[w]$
  असेल तर $\mSat{M}{!A}[w]$ असते. $\Gamma$ मध्ये
  एकच !!{formula}~$!B$ असेल, तर $!B \Entails !A$ लिहितो.
\end{defn}

\begin{ex}
  एका !!a{formula} पासून दुसरे तार्किकरीत्या निष्पन्न होते हे
  दाखवण्यासाठी $\mSat{M}{}[w]$ ची व्याख्या वापरून सर्व
  प्रतिमानांबद्दल विचार करावा लागतो. उदाहरणार्थ,
  $p \lif \Diamond p \Entails \Box\lnot p \lif \lnot p$
  दाखवण्यासाठी असा युक्तिवाद करता येईल: प्रतिमान
  $\mModel{M} = \tuple{W, R,
    V}$ आणि $w \in W$ घ्या, आणि $\mSat{M}{p \lif \Diamond
    p}[w]$ असे समजा. $\mSat{M}{\Box\lnot p \lif \lnot
    p}[w]$ दाखवायचे आहे. तसे नाही असे समजा. मग
  $\mSat{M}{\Box\lnot p}[w]$ आणि $\mSat/{M}{\lnot p}[w]$.
  $\mSat/{M}{\lnot p}[w]$ असल्यामुळे $\mSat{M}{
    p}[w]$. गृहीतकावरून $\mSat{M}{p \lif \Diamond p}[w]$;
  म्हणून $\mSat{M}{\Diamond p}[w]$.
  $\mSat{M}{\Diamond
    p}[w]$ च्या व्याख्येवरून $Rww'$ असलेले असे $w'$
  आहे की $\mSat{M}{p}[w']$. पण $\mSat{M}{\Box \lnot p}[w]$
  देखील असल्यामुळे $\mSat{M}{\lnot p}[w']$; हा विरोधाभास.

  !!a{formula}~$!B$ पासून दुसरे~$!A$ तार्किकरीत्या
  निष्पन्न होत नाही हे दाखवण्यासाठी प्रतिदृष्टांत द्यावा
  लागतो: असे प्रतिमान $\mModel{M} = \tuple{W, R,
    V}$ की त्यातील एखाद्या जगात~$w \in W$ येथे
  $\mSat{M}{!B}[w]$, पण $\mSat/{M}{!A}[w]$.
  $p \lif \Diamond p \Entails/
  \Box p \lif p$ हे दाखवू. \olref{fig:counterex} मधील
  प्रतिमान पाहा. येथे $\mSat{M}{\Diamond p}[w_1]$
  आणि म्हणून $\mSat{M}{p \lif
    \Diamond p}[w_1]$. मात्र $\mSat{M}{\Box p}[w_1]$
  असूनही $\mSat/{M}{p}[w_1]$ असल्यामुळे
  $\mSat/{M}{\Box p \lif p}[w_1]$.
  \begin{figure}
  \begin{center}
    \begin{tikzpicture}[modal]
      \node[world] (w1) [label=right:\mFalse{p}]{$w_1$};
      \node[world] (w2) [label=right:\mTrue{p}, above left=of w1]{$w_2$};
      \node[world] (w3) [label=right:\mTrue{p}, above right=of w1] {$w_3$};
      \draw[->] (w1) to (w2);
      \draw[->] (w1) to (w3);
    \end{tikzpicture}
  \end{center}
  \caption{$p \lif \Diamond p
  \Entails \Box p \lif p$ याचा प्रतिदृष्टांत.}
  \ollabel{fig:counterex}
  \end{figure}

  अनेकदा अगदी साधे प्रतिदृष्टांत पुरेसे असतात.
  $W' = \{w\}$, $R' = \emptyset$, आणि $V'(p) =
  \emptyset$ असलेले प्रतिमान $\mModel{M'} =
  \tuple{W', R', V'}$ हाही प्रतिदृष्टांत आहे:
  $\mSat/{M'}{p}[w]$ असल्यामुळे
  $\mSat{M'}{p \lif \Diamond p}[w]$. $w$ पासून एकही
  जग प्राप्य नसल्यामुळे $\mSat{M'}{\Box p}[w]$; म्हणून
  $\mSat/{M'}{\Box p
    \lif p}[w]$.
\end{ex}

\begin{prob}
  $\Box (!A \land !B) \Entails \Box !A$ हे दाखवा.
\end{prob}

\begin{prob}
  $\Box (p \lif q) \Entails/ p \lif \Box q$ आणि $p \lif \Box
  q \Entails/\Box (p \lif q)$ हे दाखवा.
\end{prob}

\end{document}
